\documentclass{article} % simplest latex class that gives nice default structures.  Other classes include standalone, book

% Recommended packages:
\usepackage{listings} % used for listing code

\usepackage{graphicx} % used for importing graphics


% Suggested packages:
\usepackage{tikz} % used for drawing figures in LaTeX directly
\usetikzlibrary{automata,calc,arrows} % some useful default tikz packages

\usepackage{enumerate} % used for finer control over enumerate lists

\usepackage{algorithm} % used for specifying algorithms
\usepackage{algpseudocode} % ditto (you need both)
\usepackage{amsmath}
\usepackage{fullpage} % reduces a page's margins to 1cm instead of 1in

\usepackage{mathpazo} % better? fonts

\usepackage{hyperref} % hyperlinks: compile with pdflatex for best results

\usepackage{amssymb}


\begin{document}

\title{Assignment 2}
\author{Tianyi Zhu} % let's keep things anonymous
\maketitle % to actually insert the above three items

\section*{Problem 1} % use section* to suppress numbering
\subsection*{(a)}
$\because R_1 \cap R_2$ has all the elements which are both in $R_1$ and $R_2$.\\
$\therefore$ If some element is both in $R_1$ and $R_2$, it is an element of $R_1 \cap R_2$.\\
$\because R_1, R_2 $ are two equivalence relations on S.\\
$\therefore R_1, R_2$ are reflexive, symmetric and transitive.\\
\\
Reflexive:\\
$\because$ According to the definition of reflexive: For all $x \in S$, $(x, x)\in R_1, R_2$.\\
$\therefore (x, x)\in R_1\cap R_2$ and $R_1 \cap R_2$ is reflexive.\\
\\
Symmetric:\\
Suppose there exists some $(x', y') \in R_1 \cap R_2$, then $(x', y') \in R_1, R_2$.\\
$\because$ According to the definition of symmetric : \\
For all $x, y \in S$: If $(x, y) \in R_1$ then $(y, x) \in R_1$. \\
For all $x, y \in S$: If $(x, y) \in R_2$ then $(y, x) \in R_2$. \\
$\because (x', y')\in R_1, R_2$\\
$\therefore (y', x') \in R_1, R_2$. Then $(y', x') \in R_1 \cap R_2$ and $R_1 \cap R_2$ is symmetric.\\
\\
Transitive:\\
Suppose there exists some $(x', y'), (y', z')\in R_1 \cap R_2$, then $(x', y'), (y', z')\in R_1, R_2$\\
$\because$ According to the definition of transitive: \\
For all $x, y, z \in S$: If $(x, y)$ and $(y, z) \in R_1$ then $(x, z) \in R_1$. \\
For all $x, y, z \in S$: If $(x, y)$ and $(y, z) \in R_2$ then $(x, z) \in R_2$. \\
$\because (x', y'), (y', z')\in R_1, R_2$\\
$\therefore (x', z') \in R_1, R_2$. Then $(x', z')\in R_1 \cap R_2$ and $R_1 \cap R_2$ is transitive.\\
\\
$\because R_1 \cap R_2$ is reflexive, symmetric and transitive. \\
$\therefore R_1 \cap R_2$ is an equivalence relation\\
\subsection*{(b)}
Suppose there exists some element $y \in [x]_1, [x]_2$\\
By the definition of $[x]_1, [x]_2$: $(x, y) \in R_1, R_2$\\
$(x, y) \in R_1 \cap R_2$ since $(x, y) \in R_1, R_2$\\
Hence all the common elements of $[x]_1$ and $[x]_2$ are in $[x]$.\\
\\
Suppose there exists element $z \in [x]$\\
By the definition of $[x]$: $(x, z) \in R_1 \cap R_2$\\
$(x, y) \in R_1, R_2$ since $(x, y) \in R_1 \cap R_2$\\
Hence all the elements of $[x]$ are in both $[x]_1$ and $[x]_2$.\\
\\
$\because$ All the common elements of $[x]_1$ and $[x]_2$ are in $[x]$, and all the elements of $[x]$ are in both $[x]_1$ and $[x]_2$.\\
$\therefore [x] = [x]_1 \cap [x]_2$\\
\subsection*{(c)}
Assume $S = \{1, 2, 3\}$, $R_1 = \{(1, 1), (2, 3), (3, 3), (1, 2), (2, 1)\}$ and $R_2 = \{(1, 1), (2, 2), (3, 3), (2, 3), (3, 2)\}$\\
Then $R_1 \cup R_2 =\{(1, 1), (2, 3), (3, 3), (1, 2), (2, 1), (2, 3), (3, 2)\}$\\
Because for $(1, 2), (2, 3)\in R_1 \cup R_2$, there is no $ (1, 3) \in R_1 \cup R_2$.\\
Hence $R_1 \cup R_2$ is not always an equivalence relation.\\

\section*{Problem 2}
\subsection*{(a)}
Because By the definition of reflexive: For all $x \in S$, $(x, x)\in R_1, R_2$\\
Therefore For all $x \in S$, $ (x, x) \in R_1;R_2$\\
Hence $R_1;R_2$ is reflexive.\\

\subsection*{(b)}
Assume that there is no common b with $(a, b)\in R_1, (b, c)\in R_2$, then $R_1;R_2$ is $\emptyset$.\\
Because $\emptyset$ is symmetric, $R_1;R_2$ is symmetric in this situation.\\
\\
Assume there exists a b with $(a, b)\in R_1, (b, c)\in R_2$.\\
By the definition of symmetric: $(b, a)\in R_1$ and $(c, b)\in R_2$\\
But there is no such b with $(c, b) \in R_1$ and $(b, a) \in R_2$\\
Therefore $(c,a) \notin R_1;R_2$\\
Hence $R_1;R_2$ is not symmetric.\\

\subsection*{(c)}
Assume $S=\{1, 2, 3\}$, $R_1=\{(2, 1), (1, 3), (2, 3)\}$ and $R_2 = \{(1, 3), (3, 2), (1, 2)\}$\\
$R_1$ is transitive because for $(2, 1), (1, 3) \in R_1$, $(2, 3) \in R_1$\\
$R_1$ is transitive because for $(1, 3), (3, 2) \in R_1$, $(1, 2) \in R_1$\\
Then $R_1;R_2 = \{(2, 3), (1, 2), (2, 2)\}$\\
Because there is $(1, 3)$ in $R_1;R_2$, $R_1;R_2$ is not transitive. 

\section*{Problem 3}
\subsection*{(a)}
Let $P(a)$ be the proposition that $R^{i+a+1}=R^{i+a}$\\
We will show that P(a) holds for all a by induction on a\\
\\
Base case: $a=0$\\
$P(0):R^i=R^{i+1}$\\
$P(0)$ holds\\
\\
Inductive case:\\
Assume P(a) hold. That is, $R^{i+a}=R^{i+a+1}$\\
We will show P(a+1) holds\\
$\because R^{i+a+1} = R^{i+a} \cup (R;R^{i+a})$\\
$\therefore(R;R^{i+a}) \subseteq R^{i+a+1}$\\
$\because R^{i+a} = R^{i+a+1}$\\
$\therefore R;R^{i+a} = R;R^{i+a+1}$\\
$\because R;R^{i+a}\subseteq R^{i+a+1}$ and $R;R^{i+a} = R;R^{i+a+1}$\\
$\therefore R;R^{i+a+1} \subseteq R^{i+a+1}$\\
$\because R;R^{i+a+1} \subseteq R^{i+a+1}$ and $R^{i+a+2} = R^{i+a+1} \cup (R;R^{i+a+1})$\\
$\therefore R^{i+a+2} = R^{i+a+1}$\\
So $P(a)$ implies $P(a+1)$.\\
Therefore, by induction, P(a) hold for all x and $R^j = R^i$ for all $j \geqslant i$.\\

\subsection*{(b)}
$\because R^j = R^i$ for all $j \geqslant i$\\
$\therefore R^j \subseteq R^i$ holds for all $j \geqslant i$\\
\\
$\because R^{n+1} = R^{n} \cup (R;R^{n})$\\
$\therefore R^n \subseteq R^{n + 1}$\\
Base case: $P(0): R^{i+1} \subseteq R^i$\\
Induction case: $P(a): R^{i-a} \subseteq R^i$\\
$\because R^n \subseteq R^{n + 1}$\\
$\therefore P(a-1): R^{i-a-1} \subseteq R^i$\\
So $P(a)$ implies $P(a-1)$.\\
Therefore for all $0\leqslant j<i$,  $R^j \subseteq R^i$ holds\\
\\
So  $R^j \subseteq R^i$ for all $j\geqslant 0$

\subsection*{(c)}
$\because R^j = R^{j+1}$ for $j\geqslant i_{min}$. \ \ ($i_{min}$ means the smallest $i$)\\
$\therefore R^{k^2} = R^{k^2+1}$ if $k^2 > i_{min}$\\
$\because R^n \subset R^{n + 1}$\\
$\therefore R^{x-1} \subset R^{x}$ for $1 < x < i_{min}$ \\
$\therefore \vert R^{x}\vert_{min} = \vert R^{x-1}\vert + 1$\\
Suppose a binary relation $(x_1, x_2) \subseteq S \times S$.\\
Both $x_1$ and $x_2$ have $k$ different value to select from. And the selection of value is independent which means the selection of $x_1$ value does not have an influence on the selection of $x_2$ value, vice versa.\\
So the max number of possible $(x_1,x_2)$ is $k^2$ which is also the maximum of $\vert R^{n}\vert$.\\
$\because \vert R^{0}\vert = k$, $\vert R^{n}\vert_{max} = k^2$ and $\vert R^{x}\vert_{min} = \vert R^{x-1}\vert + 1$\\
$\therefore i_{min}\leqslant k^2-k$\\
$\because k^2-k < k$\\
$\therefore R^{k^2} = R^{k^2+1}$ holds for $\vert S \vert = k$\\

\subsection*{(d)}
Let $P(n)$ be the proposition that for all $m \in N: R^n; R^m = R^{n+m}$\\
Base case: $P(0)$\\
\[
\begin{array}{rlc}
R^{0};R^m &= I;R^m\\
        &= R^m &\mbox{(Assignment 1 8(b))}\\
        &= R^{m+0}\\
\end{array} 
\]
So $P(0)$ holds.
\\
Inductive case:\\
Assume $P(n)$ holds that is for all $m \in N: R^n; R^m = R^{n+m}$\\
\[
\begin{array}{rlc}
    R^1;R^m     &= (R^0 \cup R;R^0);R^m \\
                &= (I \cup R;I);R^m \\
                &= (I \cup R);R^m       &\mbox{(Assignment 1 8(b))}\\
                &= (I;R^m) \cup (R;R^m)   &\mbox{(Assignment 1 8(c))}\\
                &= R^m \cup (R;R^m)     &\mbox{(Assignment 1 8(b))}\\
                &= R^{m+1}\\
\end{array}
\]
So $R^1;R^m = R^{m + 1}$ holds for $m \in \mathbb{N}$.\\
\[
\begin{array}{rlc}
    R^1;(R^n;R^m)   &= R^1;(R^{n + m})    \\
    (R^1;R^n);R^m   &= R^{n + m + 1} \\
    R^{n + 1};R^{m} &= R^{n + m + 1}\\
          
\end{array}
\]
So $P(n)$ implies $P(n + 1)$\\
Therefore, by induction, $P(n)$ holds for all $n\in \mathbb{N}$.

\subsection*{(e)}
Assume there is $(x, y), (y, z)\in R^m$ and $(x, z) \notin R^m$\\
According to (d), $R^{2m} = R^m;R^m$.\\
By the definition of ; : $(x, z) \in R^{2m}$\\
$\because R^{2m}\subseteq R^{k^2}$\\
$\therefore$ For all $x, y, z \in S$: If $(x, y)$ and $(y, z)\in R^{k^2}$, then ${x, z}\in R^{k^2}$\\
\\
Assume there no $(x, y), (y, z)\in R^m$ for all m\\
\[\begin{array}{rlcr}
    R^1 &= R^0 \cup (R;R^0)\\
        &= I \cup R\\
\end{array}\]
Because there is no $(x, y), (y, z)\in R^m$ for all m, $(x, y), (y, z) \notin R^{k^2}$\\

So $R^{k^2}$ is transitive.\\

\subsection*{(f)}
$\because R \subseteq S\times S$ \\
$\therefore R^{\gets} \subseteq S \times S$ by the definition of $R^{\gets}$\\
$\because R \subseteq S\times S$ and $R^{\gets} \subseteq S \times S$\\
$\therefore R\cup R^{\gets} \subseteq S \times S$\\
\subsubsection*{Reflexive}
$\because$ According to (b)(c), $R^0 \subseteq (R\cup R^{\gets})^{k^2}$\\
$\because R^0 = \{(x, x) : x \in S\}$ \\
$\therefore$ For all $x\in S$, $(x, x)\in (R\cup R^{\gets})^{k^2}$\\
So $(R\cup R^{\gets})^{k^2}$ is reflexive.\\
\subsubsection*{Symmetric}
$\because R^{\gets}=\{(t,s) \in S \times S : (s,t) \in R\}$\\
$\therefore$ For all $(x, y)\in R$, then  $(y, x)\in R^{\gets}$\\
$\because (R \cup R^{\gets})^1 = I \cup (R \cup R^{\gets};I) =I \cup (R \cup R^{\gets})$\\
$\because$ According to (b)(c), $R^1 \subseteq (R\cup R^{\gets})^{k^2}$\\
$\therefore$ For all $(x, y)\in (R\cup R^{\gets})^{k^2}$, then  $(y, x)\in (R\cup R^{\gets})^{k^2}$\\
So $(R\cup R^{\gets})^{k^2}$ is symmetric.
\subsubsection*{Transitive}
$\because$ According to (e), $R^{k^2}$ is transitive for all $R \subseteq S\times S$.\\
$\therefore (R\cup R^{\gets})^{k^2}$ is transitive.\\
\\
So $(R\cup R^{\gets})^{k^2}$ is an equivalence relation.\\
\subsection*{(g)}
Assume that $(R\cup R^{\gets})^{k^2}$ is not the minimun of all the equivalence relation that contains $R$.\\
There exists some $(x, y) \notin R\cup R^{\gets}$ that can be removed from $(R\cup R^{\gets})^{k^2}$ and $(R\cup R^{\gets})^{k^2}\backslash (x, y)$ is still an equivalence relation.\\
Because $(R\cup R^{\gets})^{k^2}\backslash (x, y)$ is reflexive, $x \neq y$.\\
Because $(R\cup R^{\gets})^{k^2}\backslash (x, y)$ is symmetric, $(y, x)$ must be removed.\\
Assume there is no $(y, z), (z, x) \in (R\cup R^{\gets})^{k^2}$\\
Because there is no $(y, z), (z, x) \in (R\cup R^{\gets})^{k^2}$, there is no $(y, z), (z, x) \in$ any $(R\cup R^{\gets})^{m}$.\\
Because $(x, y)$ is impossible to get from $(R\cup R^{\gets})^n;(R\cup R^{\gets})^m$, $(x, y)$ must be in $(R\cup R^{\gets})^{1}$.\\
Because $(R\cup R^{\gets})^1 = (R\cup R^{\gets}) \cup I$ and $(x, y) \notin I$, $(x, y)$ must be in $R\cup R^{\gets}$\\
So there is a conflict and $(R\cup R^{\gets})^{k^2}$ is the minimun of all the equivalence relation that contains $R$ in this situation.\\
\\
Assume there is $(y, z), (z, x) \in (R\cup R^{\gets})^{k^2}$\\
Because $(x, y)$ is removed, $(y, z), (z, x), (z, y), (x, z)$ must also be reomved.\\
Because $(x, y), (y, x), (x, z), (z, x), (y, z), (z, y)$ are removed and none of them is an element of $R\cup R^{\gets}$\\
Because $(x, y), (y, x), (x, z), (z, x), (y, z), (z, y) \notin I$, none of them is in $(R\cup R^{\gets})^{1}$\\
So $(x, y), (y, x), (x, z), (z, x), (y, z), (z, y)$ is impossible to get from $(R\cup R^{\gets})^n;(R\cup R^{\gets})^m$.\\
$(x, y), (y, x), (x, z), (z, x), (y, z), (z, y)\notin (R\cup R^{\gets})^{k^2}$\\
There is a conflict.\\
\\
Hence, $(R\cup R^{\gets})^{k^2}$ is the minimun of all the equivalence relation that contains $R$.\\



\section*{Problem 4}
Assume $f(n) \leqslant cn$ for all positive numbers less than $n$. Therefore\\
\[
\begin{array}{rlcr}
f(n)    &\leqslant f(\lfloor \frac{n}{3} \rfloor) + 3f(\lfloor \frac{n}{5} \rfloor) +n\\
        &\leqslant c \frac{n}{3} + 3c\frac{n}{5} + n\\
        &\leqslant (\frac{14}{15}c + 1)n\\
        &\leqslant c_{next}n
\end{array}
\]
\\
$\because \frac{14}{15}<1$ is a negative feedback which means $c_{next}$ will decrease after some point.\\
$\therefore$ There is a upper bound of $c_{next}$.\\
$\because c_{next} = \frac{14}{15}c =\frac{f(n)}{n}$\\
$\therefore c_{x} =(\frac{14}{15})^x c_0 + 15(1-(\frac{14}{15})^{x})$\\
$\therefore \lim\limits_{x\to+\infty}c_{x} =(\frac{14}{15})^x c_0 + 15(1-(\frac{14}{15})^{x}) = 15$\\
So if $f(n) \leqslant cn$ for all positive numbers less than $n$, $f(n)\leqslant cn$\\
$\because$
\[
\begin{array}{rlcr}
    f(1)    &=f(0)+3f(0)+1=1 \\
    f(2)    &=f(0)+3f(0)+2=2 \\
\end{array}
\]
$\therefore f(n)\leqslant n$ for all positive numbers less than 3\\
So $f(n) \in O(n)$\\

\section*{Problem 5}
\subsection*{(a)}
count(T):\\
.\qquad if $T == \tau$\\
.\qquad\qquad return 0\\
.\qquad elif $T_{left} == \tau$ and $T_{right} == \tau$\\
.\qquad\qquad return 1\\
.\qquad elif $T_{left} == \tau$\\
.\qquad\qquad return 1 + count($T_{right}$)\\
.\qquad elif $T_{right} == \tau$\\
.\qquad\qquad return 1 + count($T_{left}$)\\
.\qquad else\\
.\qquad\qquad return 1 + count($T_{left}$) + count($T_{right}$)\\
\\
\subsection*{(b)}
leaves(T):\\
.\qquad if $T == \tau$\\
.\qquad\qquad return 0\\
.\qquad elif $T_{left} == \tau$ and $T_{right} == \tau$\\
.\qquad\qquad return 1\\
.\qquad elif $T_{left} == \tau$\\
.\qquad\qquad return leaves($T_{right}$)\\
.\qquad elif $T_{right} == \tau$\\
.\qquad\qquad return leaves($T_{left}$)\\
.\qquad else\\
.\qquad\qquad return leaves($T_{left}$) + leaves($T_{right}$)\\

\subsection*{(c)}
half-leaves(T):\\
.\qquad if $T == \tau$\\
.\qquad\qquad return 1\\
.\qquad elif $T_{left} == \tau$ and $T_{right} == \tau$\\
.\qquad\qquad return 0\\
.\qquad elif $T_{left} == \tau$\\
.\qquad\qquad return 1 + half-leaves($T_{left}$)\\
.\qquad elif $T_{right} == \tau$\\
.\qquad\qquad return 1 + half-leaves($T_{right}$)\\
.\qquad else\\
.\qquad\qquad return half-leaves($T_{left}$) + half-leaves($T_{right}$)\\


\subsection*{(d)}
Let $P(T)$ be the proposition that $count(T) = 2 \times leaves(T) + half-leaves(T) - 1$\\
\subsubsection*{Base case: $P(\tau)$ and $P(\tau,\tau)$}
\subsubsection*{$P(\tau)$}
An empty tree has 0 node, 0 leaf and 1 half-leaf.\\
$count(\tau) = 0$\\
$2\times leaves((\tau)+half$-$leaves((\tau) - 1 = 2\times 0 + 1-1 = 0$\\
$\because count((\tau,\tau)) = 2\times leaves((\tau,\tau))+half$-$leaves(((\tau,\tau))$\\
$\therefore P((\tau)$ holds\\

\subsubsection*{$P((\tau,\tau))$}
This kind of binary tree has 1 node, 1 leaf and 0 half-leaf.\\
$count((\tau,\tau)) = 1$\\
$2\times leaves(((\tau,\tau))+half$-$leaves((\tau,\tau)) - 1 = 2\times 1 + 0-1 = 1$\\
$\because count((\tau,\tau)) = 2\times leaves((\tau,\tau))+half$-$leaves(((\tau,\tau))$\\
$\therefore P(((\tau,\tau))$ holds\\


\subsubsection*{Inductive case:}
To extend a existed binary, there are two ways. The first one extends from a leaf while the second one extends from a half-leaf\\
Assume T has N nodes, L leaves and H half-leaves. \\
\subsubsection*{Transform one leaf in T from $(\tau, \tau )$ into $((\tau,\tau),\tau)$}
Suppose the transformed binary tree is called $T_{new1}$\\
$count(((\tau,\tau),\tau))$ = 1 + $count(\tau, \tau)$\\
$leaves(((\tau,\tau),\tau)) = leaves((\tau,\tau))$\\
$half-leaves(((\tau,\tau),\tau)) = 1 + half-leaves((\tau,\tau))$\\
\\
$N_{new1} = N + 1$\\
$L_{new1} = L $\\
$H_{new1} = H + 1$\\
\\
$\because N = 2\times L + H - 1$\\
$\therefore N_{new1} = 2\times L_{new1} + H_{new1} - 1$\\
So $P(T)$ implies $P(T_{new1})$\\
\subsubsection*{Transform one leaf in T from $(\tau, \tau )$ into $(\tau,(\tau,\tau))$}
Suppose the transformed binary tree is called $T_{new2}$\\
$count((\tau,(\tau,\tau)))$ = 1 + $count(\tau, \tau)$\\
$leaves((\tau,(\tau,\tau))) = leaves((\tau,\tau)$\\
$half-leaves((\tau,(\tau,\tau))) = 1 + half-leaves((\tau,\tau))$\\
\\
$N_{new2} = N + 1$\\
$L_{new2} = L $\\
$H_{new2} = H + 1$\\
\\
$\because N = 2\times L + H - 1$\\
$\therefore N_{new2} = 2\times L_{new2} + H_{new2} - 1$\\
So $P(T)$ implies $P(T_{new2})$\\

\subsubsection*{Transform one half-leaf in T from $(t_{left}, \tau )$ into $(t_{left},(\tau,\tau))$}
Suppose the transformed binary tree is called $T_{new3}$\\
$count((t_{left},\tau)) = count(t_{left}) + 1$\\
$count((t_{left},(\tau,\tau))) = count(t_{left}) + count((\tau,\tau))+1 = count(t_{left}) + 2$\\
$leaves((t_{left},\tau)) = leaves(t_{left})$\\
$leaves((t_{left},(\tau,\tau)) = leaves(t_{left}) + leaves((\tau,\tau)) = leaves(t_{left}) + 1$\\
$half-leaves((t_{left},\tau)) = leaves(t_{left}) + 1$\\
$half-leaves((t_{left},(\tau,\tau)) = half-leaves(t_{left}) + half-leaves(\tau, \tau) = half-leaves(t_{left})$\\
\\
$N_{new3} = N + 1$\\
$L_{new3} = L + 1$\\
$H_{new3} = H - 1$\\
\\
$\because N = 2\times L + H - 1$\\
$\therefore N_{new3} = 2\times L_{new3} + H_{new3} - 1$\\
So $P(T)$ implies $P(T_{new3})$\\
\subsubsection*{Transform one half-leaf in T from $(\tau, t_{right})$ into $((\tau,\tau), t_{right})$}
Suppose the transformed binary tree is called $T_{new4}$\\
$count((\tau, t_{right})) = count(t_{right}) + 1$\\
$count((\tau,\tau), t_{right}) = count(t_{right}) + count((\tau,\tau))+1 = count(t_{right}) + 2$\\
$leaves((\tau, t_{right})) = leaves(t_{right})$\\
$leaves((\tau,\tau), t_{right}) = leaves(t_{right}) + leaves((\tau,\tau)) = leaves(t_{right}) + 1$\\
$half-leaves((\tau, t_{right})) = leaves(t_{right}) + 1$\\
$half-leaves((\tau,\tau), t_{right}) = half-leaves(t_{right}) + half-leaves(\tau, \tau) = half-leaves(t_{right})$\\
\\
$N_{new4} = N + 1$\\
$L_{new4} = L + 1$\\
$H_{new4} = H - 1$\\
\\
$\because N = 2\times L + H - 1$\\
$\therefore N_{new4} = 2\times L_{new4} + H_{new4} - 1$\\
So $P(T)$ implies $P(T_{new4})$\\
\\
Therefore, by induction, P(T) holds for all T.\\

\section*{Problem 6}
\subsection*{(a)}
The running time of sum is $O(1)+O(n)+O(n^2)+O(1)$ and its asymptotic upper bound is $O(n^2)$\\
\subsection*{(b)}
The running time of product is $O(1)+O(n)+O(n^4)+O(1)$ and its asymptotic upper bound is $O(n^4)$\\
\subsection*{(c)}
One multiplication of two $n \times n$ matrices includes 8 multiplications of two $\frac{n}{2} \times \frac{n}{2}$ matrices 
and 4 additions of two $\frac{n}{2} \times \frac{n}{2}$ matrices.\\
The running time of one multiplications of two $\frac{n}{2} \times \frac{n}{2}$ matrices is $T(\frac{n}{2})$\\
The running time of one additions of two $\frac{n}{2} \times \frac{n}{2}$ matrices is $(\frac{n}{2})^2$\\
Therefore, the recurrence equation of $T(n)$ is $T(n) = 8T(\frac{n}{2}) +4(\frac{n}{2})^2$

\subsection*{(d)}
$T(n) = 8T(\frac{n}{2}) +4(\frac{n}{2})^2$\\
According to Master Theorem, $a=8$, $b=2$, $c=2$ and $d=\log_{b}{a}=3$\\
$\because d<c$\\
$\therefore T(n)=O(n^3)$\\


\end{document}

